\documentclass[UTF8,12pt]{ctexart}
\usepackage[a4paper,margin=24mm]{geometry}
\usepackage{amsmath,amssymb,bm,graphicx,adjustbox}
\newcommand{\fitmath}[1]{\begin{adjustbox}{max width=\linewidth}$\displaystyle #1$\end{adjustbox}}
\setlength{\parindent}{0pt}
\setlength{\parskip}{.7em}
\sloppy
\allowdisplaybreaks
\begin{document}
\section*{1999 年考研数学三 · 参考答案}
\subsection*{1999 年真题参考答案}

\subsection*{一、填空题}

（1）\(\displaystyle \dfrac{\displaystyle 4}{\displaystyle \pi}-1\)。（2）\(\displaystyle 4\)。（3）\(\displaystyle O_{3\times3}\)，即 3 阶零矩阵。（4）\(\displaystyle 16\)。（5）\(\displaystyle 0\)。


\subsection*{二、选择题}

（1）A。（2）C。（3）B。（4）D。（5）A。


三、所求切线方程为 \(\displaystyle y-\dfrac{\displaystyle 1}{\displaystyle \sqrt a}=-\dfrac{\displaystyle 1}{\displaystyle 2\sqrt{a^3}}(x-a)\)，面积为 \(\displaystyle \dfrac{\displaystyle 9}{\displaystyle 4}\sqrt a\)。当切点沿 \(\displaystyle x\) 轴正方向趋于无穷远时，\(\displaystyle \lim_{a\to+\infty}S=+\infty\)；当切点沿 \(\displaystyle y\) 轴正方向趋于无穷远时，\(\displaystyle \lim_{a\to0^+}S=0\)。


四、\(\displaystyle 4-\dfrac{\displaystyle \pi}{\displaystyle 2}\)。


五、当 \(\displaystyle x_1=6\left(\dfrac{\displaystyle p_2\alpha}{\displaystyle p_1\beta}\right)^\beta\)，\(\displaystyle x_2=6\left(\dfrac{\displaystyle p_1\beta}{\displaystyle p_2\alpha}\right)^\alpha\) 时，投入总费用最小。


六、\(\displaystyle y(x)=\begin{cases}e^{2x}-1,\allowbreak &x\le1,\allowbreak \\(1-e^{-2})e^{2x},\allowbreak &x>1.\end{cases}\)


七、\(\displaystyle \dfrac{\displaystyle 3}{\displaystyle 4}\)。


八、（1）证明略。（可利用零点定理。）（2）证明略。（可考虑函数 \(\displaystyle F(x)=e^{-\lambda x}[f(x)-x]\)，对 \(\displaystyle F(x)\) 使用罗尔定理。）


九、\(\displaystyle a=2,\allowbreak b=-3,\allowbreak c=2,\allowbreak \lambda_0=1\)。


十、证明略。（\(\displaystyle B\) 为对称矩阵，\(\displaystyle x^{\mathrm T}Bx=\lambda x^{\mathrm T}x+(Ax)^{\mathrm T}(Ax)\)。）


十一、（1）\(\displaystyle P\{U=0,\allowbreak V=0\}=\dfrac{\displaystyle 1}{\displaystyle 4}\)，\(\displaystyle P\{U=0,\allowbreak V=1\}=0\)，\(\displaystyle P\{U=1,\allowbreak V=0\}=\dfrac{\displaystyle 1}{\displaystyle 4}\)，\(\displaystyle P\{U=1,\allowbreak V=1\}=\dfrac{\displaystyle 1}{\displaystyle 2}\)。（2）\(\displaystyle r=\dfrac{\displaystyle 1}{\displaystyle \sqrt3}\)。


十二、证明略。

\end{document}
